\section{Conclusion}

Deep-Dream presents a source-grounded memory design in which structured concepts and relations help an agent navigate long histories without replacing the histories that authorize an answer. The current implementation combines document-first ingestion, conservative alignment, hybrid retrieval, source-neighbor expansion, and a scope-checked evidence interface. Full LoCoMo-1540 and LongMemEval-S-500 runs show strong protocol-aligned internal results, while retrieval diagnostics make the cost of wider evidence contexts visible.

The paper's main scientific boundary is also its main limitation. Valid evidence identifiers do not guarantee that a generated answer is semantically entailed; LongMemEval update scores do not by themselves prove that versioning caused the improvement; and MEME deletion at 0\% shows that negative-memory behavior remains unsolved. We therefore separate completed evidence from the required next experiments: causal provenance/evidence-gate ablations, update and deletion tests, repaired budget curves, and controlled Mem0/Graphiti comparisons. These experiments will determine whether the auditable overlay is a useful mechanism beyond a carefully instrumented system design.
